\section{Related Work}

\subsection{Large Language Models on Medical Benchmarks}

The evaluation of LLMs on medical knowledge benchmarks has grown into a substantial research area since the release of GPT-3 \parencite{brown2020gpt3}. The MedQA dataset, introduced by \textcite{jin2021medqa}, established a widely-used benchmark of USMLE-style multiple-choice questions derived from publicly released practice materials, enabling reproducible comparisons across models. Subsequent benchmarks—including MedMCQA \parencite{pal2022medmcqa}, which covers medical entrance examination questions from India, and PubMedQA \parencite{jin2019pubmedqa}, which tests biomedical literature comprehension—have expanded the evaluation landscape beyond the USMLE context.

The landmark study by \textcite{nori2023gpt4usmle} demonstrated that GPT-4 achieved accuracy exceeding the passing threshold on all three USMLE steps, with an overall score of approximately 87\% on Step 1 and Step 2 CK questions—a substantial improvement over GPT-3.5, which had been shown by \textcite{kung2023chatgptusmle} to perform near the passing threshold in an earlier study. \textcite{singhal2023medpalm} introduced Med-PaLM, a domain-adapted variant of PaLM fine-tuned on medical QA pairs, achieving 67.2\% on the MedQA benchmark; its successor Med-PaLM 2 reached over 85\% \parencite{singhal2023medpalm2}. More recently, OpenAI's o1 reasoning model demonstrated state-of-the-art performance on complex clinical vignettes, suggesting that extended chain-of-thought reasoning provides substantial gains on medical MCQs \parencite{openai2024o1}.

However, these studies share a common limitation: they evaluate one or two models in isolation, making it difficult to situate any single model's performance within the broader landscape of currently available systems. Our study addresses this by evaluating eight models in a unified, controlled setting. Additionally, most prior work does not disaggregate results by exam step or medical subject in sufficient granularity to support actionable guidance for students preparing for specific examination components.

\subsection{AI Tools in Medical Education}

The adoption of AI chatbots by medical students has been rapid and widespread. Surveys conducted across US and European medical schools in 2023 found that between 60\% and 80\% of students reported using ChatGPT or similar tools for course preparation, with the most common use cases being concept explanation, case-based reasoning practice, and question-answer review \parencite{huh2023medstudents, karabacak2023chatgptmed}. Several institutions have begun formally integrating AI tools into curricula, while others have raised concerns about accuracy and hallucination risks in clinical contexts \parencite{hartsfield2023chatgptmed}.

\textcite{kanjee2023chatgptusmle} evaluated ChatGPT on a sample of USMLE questions and found performance comparable to a third-year medical student, with notable weaknesses in questions requiring multi-step reasoning or synthesis of laboratory data. Critically, the authors also documented cases of confident, plausible-sounding but factually incorrect reasoning—a pattern sometimes called ``hallucination with authority'' that poses particular risks in a study context where students may not have sufficient domain knowledge to detect errors.

\textcite{chen2023llmmededucation} conducted a systematic review of 35 studies evaluating LLMs in medical education contexts and found that performance varied substantially by domain, question type, and evaluation methodology, underscoring the need for standardized, reproducible benchmarks. Our study contributes to this standardization effort by fixing the random seed for question sampling, using temperature=0 for all model calls, and releasing all code and data to support replication.

\subsection{IMG-Specific Challenges on the USMLE}

International medical graduates represent a heterogeneous group that includes US citizens who attended medical school abroad, as well as foreign nationals who seek licensure and residency training in the United States. Despite comprising a large fraction of the physician workforce—particularly in underserved specialties and geographic areas—IMGs consistently show lower first-attempt pass rates on all USMLE steps compared to graduates of accredited US and Canadian programs \parencite{ecfmg2023annual}.

Research into the sources of this disparity has identified several factors beyond academic preparation. \textcite{boulet2006} found that questions requiring knowledge of US healthcare system operations, insurance reimbursement patterns, and regulatory frameworks posed particular difficulties for non-US graduates. Similarly, \textcite{seelig2010img} documented that clinical vignettes anchored in North American epidemiological baselines—disease prevalence patterns, demographic risk factors, and environmental exposures specific to the US population—systematically disadvantaged candidates whose training occurred in different epidemiological contexts.

Despite the growing use of LLMs among IMG students as accessible, low-cost study tools, no prior study has examined whether these models themselves encode the same US-centric biases found in the examination, or whether differential accuracy on US-centric questions might lead IMG students to receive less reliable guidance on precisely the question types where they face the greatest challenge. Our IMG perspective analysis module addresses this gap directly, using a validated keyword lexicon to identify US-centric questions and quantifying accuracy differentials using Fisher's exact test.
